\scanpage{110}\mspage{5}
La donnée d'un point $x$ \inserted{$\ne0$} de $E(k)$ équivaut à celle de son image dans $\mathbb P^1_k$, soit \deleted{$p(x)$}
\[
 p(x)=x_0\in k\,, \tag*{(38)}
\]
plus la donnée d'une des deux solutions \inserted{$x'$} de l'équation (\uline{deux} tout au moins si le $2^{\mathrm e}$ mb est $\ne0$ i.e. si $x\notin J={}_2E^*$ i.e. $x_0\notin D_J\subset\Sigma_J$)
\[
 {x'}^2=x_0^3+a_1x_0^2+a_2x_0+a_3\,, \tag*{(39)}
\]
donnée en termes de $x$ par
\[
 x'=p'(x)\,. \tag*{(40)}
\]
\uline{Proposition} 3 Soient $x,y,z\in E(k)\smallsetminus\{0\}$, \inserted{\deleted{pas} tous \deleted{égaux}\inserted{\uline{distincts}},} alors on a
\[
 x+y+z=0 \tag*{(\deleted{\unsure{4}}41)}
\]
SSS on a \deleted{a) On \unsure{n'a pas|a} $p(x)\unsure{=|\ne}p(y)=p(z)$ et b) on a}
\[
 \det\begin{pmatrix}
 1&p(x)&p'(x)\\
 1&p(y)&p'(y)\\
 1&p(z)&p'(z)
 \end{pmatrix}=0 \tag*{(42)}
\]
% slanted
\margin{\uline{NB} Si \deleted{deux} parmi les pts $x,y,z$ sont égaux, (42) est tjrs vérifié, sans qu'on ait \unsure{néc$^{\mathrm t}$} pour autant $x+y+z=0$ !}
% slanted
\margin{Cela implique qu'on a $p(x)\ne p(y)\ne p(z)\ne p(x)$ (car si on avait $p(x)=p(y)$ on aurait $x=y$ ou $x=-y$ d'où $z=0$), l'un et l'autre exclus}
% slanted
\margin{\deleted{La condition \unsure{a)|b)} \unsure{est vraie|peut même}, \unsure{car non a)|cas non a)} \unsure{signifierait|signifier} $\exists a\in E(k)$, avec deux \deleted{\illegible{1}} \deleted{\unsure{autres}} parmi les trois $x,y,z$ égaux à $a$, le troisième à $-a$, d'où $x+y+z=\unsure{a+a-a=a|a,\ \text{et}\ -a=a}$, donc $a=0$, ce qui est absurde puisque $x,y,z\ne0$.}}
En effet, la définition de la loi de multiplication de $E$ nous donne que si $x,y,z\in E(k)$ \unsure{autre condition|autre convention}, les \deleted{$x+y+z=0$}
\[
 \begin{aligned}
 x+y+z=0&\iff(\{x\}-\{0\})+(\{y\}-\{0\})+(\{z\}-\{0\})\overset{\rm lin.}{\sim}0\\
 &\text{i.e.}\quad\{x\}+\{y\}+\{z\}\overset{\rm lin.}{\sim}3\{0\}
 \end{aligned}
\]
Si $x,y,z$ sont tous $\ne0$, cela signifie qu'il existe une fonction rat. sur $E$, ayant $0$ comme pôle d'ordre 3, pas d'autre pôle, (donc $f\in F_3$), et $\{x\}+\{y\}+\{z\}$ comme \deleted{pôle} diviseur des zéros.
% archive sheet number: 109
